\documentclass[a4paper,11pt]{ltjsarticle}
\usepackage[haranoaji]{luatexja-preset}
\usepackage{amsmath,amssymb}
\usepackage[top=12mm,bottom=10mm,left=14mm,right=14mm]{geometry}
\usepackage{array}
\usepackage{tikz}
\usetikzlibrary{calc}
\pagestyle{empty}
\setlength{\parindent}{0pt}
\newlength{\qcol}\setlength{\qcol}{0.47\textwidth}
\newlength{\qgap}
\newlength{\anscolw}\setlength{\anscolw}{0.30\textwidth}
\newlength{\ansh}
\newcommand{\Q}[2]{\parbox[t]{\qcol}{\raggedright (#1)\hspace{0.6em}$\displaystyle #2$}}
\newcommand{\QR}[4]{\Q{#1}{#2}\hfill\Q{#3}{#4}\par\vspace{\qgap}}
\newcommand{\anscell}[1]{\rule[-\ansdrop]{0pt}{\ansh}(#1)}
\newlength{\ansdrop}

\begin{document}
{\large\bfseries 計算問題　1}\hfill 名前(\underline{\hspace{32mm}})\par\vspace{3mm}
■（1）〜（16）の計算をしなさい。（17）〜（20）は方程式を解きなさい。\par\vspace{3mm}
\setlength{\qgap}{5.4mm}
\QR{1}{-7+(-5)}{2}{(-6)-(+9)}
\QR{3}{(2a+6)+(3a-2)}{4}{(5x+2)-(2x-5)}
\QR{5}{(-7)\times(-4)}{6}{(-36)\div(+12)}
\QR{7}{5a\times(-2)}{8}{-32x\div4}
\QR{9}{(+8)+(-6)+(-1)}{10}{6\div15\times(-20)}
\QR{11}{\dfrac{x-1}{2}\times4}{12}{(12a-6)\div2}
\QR{13}{4+3^2+(-5)}{14}{15-(-21)\div\left(-\dfrac{3}{4}\right)}
\QR{15}{\dfrac{1}{2}(4x+2)+\dfrac{1}{3}(6x+3)}{16}{\dfrac{1}{4}(8a+4)-\dfrac{1}{2}(2a-6)}
\QR{17}{3x+1=-8}{18}{4(x+3)=3x+7}
\QR{19}{0.3x-0.5=1}{20}{\dfrac{1}{2}x=\dfrac{1}{3}x+2}
\vspace{1mm}
\Q{21}{12^2}\par\vspace{3mm}
\setlength{\ansh}{9.5mm}\setlength{\ansdrop}{6mm}%
\begin{center}\begin{tabular}{|p{\anscolw}|p{\anscolw}|p{\anscolw}|}\hline
\anscell{1} & \anscell{2} & \anscell{3} \\\hline
\anscell{4} & \anscell{5} & \anscell{6} \\\hline
\anscell{7} & \anscell{8} & \anscell{9} \\\hline
\anscell{10} & \anscell{11} & \anscell{12} \\\hline
\anscell{13} & \anscell{14} & \anscell{15} \\\hline
\anscell{16} & \anscell{17} & \anscell{18} \\\hline
\anscell{19} & \anscell{20} & \anscell{21} \\\hline
\end{tabular}\end{center}

\end{document}
